\documentclass[a4paper,12pt,oneside]{book}
\usepackage[czech]{babel}
\usepackage[utf8]{inputenc}

% Plugin na sloupce
\usepackage{multicol}

% Plugin na obrazky
\usepackage{graphicx}
\graphicspath{ {img} }

% Toto je plugin diky kteremu pujde klikat do obsahu
\usepackage{hyperref}
\hypersetup{
    colorlinks,
    citecolor=black,
    filecolor=black,
    linkcolor=black,
    urlcolor=black
}

% Plugin na hezci nadpis kapitoly
% https://texblog.org/2012/07/03/fancy-latex-chapter-styles/
\usepackage[T1]{fontenc}
\usepackage{titlesec, blindtext, color}
\definecolor{gray75}{gray}{0.75}
\newcommand{\hsp}{\hspace{20pt}}
\titleformat{\chapter}[hang]{\Huge\bfseries}{\thechapter\hsp\textcolor{gray75}{|}\hsp}{0pt}{\Huge\bfseries}


% Zakladni info o dokumento
\title{Ročníková práce}
\author{Lukáš Dulík}
\date{\today} % Tady sa pak moze nastavit jine datum, ted to da vzdycky aktualni


\usepackage{fancyhdr}


\fancyhf{}
\lhead{Genetika a genetické inženýrství}
\lfoot{Ročníková práce KVARTA}
\rfoot{\thepage}

\begin{document}

% Zrusi cislovani na zacatecnich strankach
\pagenumbering{gobble}

\begin{titlepage}
    \begin{center}
        \vspace*{1cm}

        \Huge
		Gymnázium Jana Pivečky a Střední odborná škola Slavičín \\

		% Logo GJP
		\includegraphics[width=0.4\textwidth]{gjp.png}

        \textbf{Ročníková práce}

        \vspace{0.5cm}
        \LARGE
        Téma: Genetika a genetické inženýrství
    \end{center}

	\vspace{1.5cm}

	% Vyplni zbyvaji prostor
	\vfill

	\vspace{0.8cm}

	\vspace{5pt}
	% Cara nahore
	\hrule
	\vspace{6pt}

	\Large

	\makeatletter
	\begin{multicols}{2}
		\noindent
		Slavičín \\
		Datum: \@date

	\columnbreak
		\noindent
		\null\hfill Třída: Kvarta \\
		\null\hfill \@author
	\end{multicols}

	\makeatother

	\vspace{5pt}

	% Cara dole
	\hrule

\end{titlepage}

" system clipboard
nmap <c-c> "+y
vmap <c-c> "+y
nmap <c-v> "+p
inoremap <c-v> <c-r>+
cnoremap <c-v> <c-r>+

\newpage
\thispagestyle{empty}
\mbox{}
\newpage


\noindent
Prohlašuji, že jsem ročníkovou práci vypracoval(a) samostatně a výhradně s použitím citovaných pramenů.

Ve Slavičíně: dne datum jméno, vlastnoruční podpis


\newpage
\section*{Abstrakt}
Abstract goes here

\newpage
\section*{Poděkování}
Děkuji paní x za pomoc s prací.

% Udela obsah
\tableofcontents

% Zacatek cislovani stranek
\pagenumbering{arabic}
\thispagestyle{fancy}

\chapter{Úvod}
\thispagestyle{fancy}

This is the first section.

Lorem  ipsum  dolor  sit  amet,  consectetuer  adipiscing
elit.   Etiam  lobortisfacilisis sem.  Nullam nec mi et
neque pharetra sollicitudin.  Praesent imperdietmi nec ante.
Donec ullamcorper, felis non sodales...


Lorem ipsum dolor sit amet, consectetuer adipiscing elit.
Etiam lobortis facilisissem.  Nullam nec mi et neque pharetra
sollicitudin.  Praesent imperdiet mi necante...

\section{Nadpis}

Lorem ipsum dolor sit amet, consectetuer adipiscing elit.
Etiam lobortis facilisissem.  Nullam nec mi et neque pharetra
sollicitudin.  Praesent imperdiet mi necante...

\subsection{Podnadpis}

Tady je neco k tomu

\subsection{Podnadpis}

Tady je neco k tomu

\section{Nadpis}

Lorem ipsum dolor sit amet, consectetuer adipiscing elit.
Etiam lobortis facilisissem.  Nullam nec mi et neque pharetra
sollicitudin.  Praesent imperdiet mi necante...

\chapter{Závěr}
\thispagestyle{fancy}

Jo jo bla bla shrnutí

\chapter{Seznam použité literatury}
\thispagestyle{fancy}

Literatura: Deník malého poseroutky xd
Literatura: Netflix Unnatural Selection

\chapter{Seznam obrázků}
\thispagestyle{fancy}

tvoja mama, s kulama

\chapter{Seznam příloh}
\thispagestyle{fancy}

Ema-Zeriavova-sextape.mp4

\end{document}
\thispagestyle{fancy}

